\documentclass{article}
\title{Practice Problem 2.49}
\author{CSSAPP}
\usepackage{booktabs}


\begin{document}
    \maketitle
    \begin{flushleft}
      A. For a floating-point format with an $n-bit$ fraction, give a formula for the 
      smallest positive integer that cannot be represented exactly (because it would require an ($n + 1$)-bit freaction to be exact).
      Assume the exponent field size $k$ is large enought that the range of representation exponent
      does not provide a limitation for this Problem.

      B. What is the numeric value of this integer for single-precision format $(n = 23)$
    
    \bigskip 
    \textbf{Solution}\\
    A. The number has a binary representation of $1$, followed by $n$ zeros, followed by $1$.
      giving the value $2^{n+1} + 1$\\

    \bigskip

    B. When $n = 23$, the value is $2^{24+1} + 1 = 16,777,217$
    
    \end{flushleft}
\end{document}